\documentclass[leqno]{article}
% Shared QBT/Soln pack preamble for resources.danieljsmith.org.
\usepackage{graphicx}
\usepackage{dsfont}
\usepackage{mathtools}
\usepackage{amssymb}
\usepackage{amsmath}
\usepackage{mleftright}
\usepackage{amsthm}
\usepackage{float}
\usepackage{ragged2e}
\usepackage{tensor}
\usepackage{fancyhdr}
\usepackage{empheq}
\usepackage[most]{tcolorbox}
\usepackage{enumitem}
\usepackage{dirtytalk}
\usepackage{marginnote}
\usepackage{microtype}
\usepackage[
  a4paper,
  left=0.8in,
  right=1in,
  top=1in,
  bottom=0.5in,
  headheight=14pt,
  headsep=0.4in
]{geometry}
\setlength{\parindent}{0cm}

% TikZ / PGFPlots
\usepackage{pgfplots}
\pgfplotsset{compat=1.18}
\usepgfplotslibrary{fillbetween, polar}
\usetikzlibrary{calc, positioning}
\usetikzlibrary{angles, quotes}
\usetikzlibrary{arrows.meta}
\usetikzlibrary{decorations.pathreplacing, decorations.markings}
\usetikzlibrary{patterns}
\usetikzlibrary{intersections}

% Links and contents. The contents follow the TMUA papers: compact entries,
% no dot leaders, and each question's mark total in grey after its label.
\usepackage{tocloft}
\tocloftpagestyle{fancy}
\setlength{\cftbeforesubsecskip}{2pt}
\setlength{\cftsubsecindent}{1.5em}
\setlength{\cftsubsecnumwidth}{0pt}
\renewcommand{\cftsubsecfont}{\small\color{black}}
\renewcommand{\cftsubsecpagefont}{\small\color{black}}
\renewcommand{\cftsubsecleader}{\hfill}
\usepackage[colorlinks=true,linkcolor=black,citecolor=magenta,urlcolor=black]{hyperref}%
\urlstyle{same}

% Pack macros
% \marks also adds its value to the current question's total in the .aux
% file, so the contents show each total from the second pass. The write
% happens at shipout, so a box typeset more than once counts once. The
% pack's mark total for the contents-page title is summed the same way.
\makeatletter
\newcount\djs@totalmarks
\newcommand{\djsAddMarks}[2]{%
  \global\advance\djs@totalmarks#2\relax
  \@ifundefined{djs@marks@#1}%
    {\expandafter\xdef\csname djs@marks@#1\endcsname{#2}}%
    {\expandafter\xdef\csname djs@marks@#1\endcsname
      {\the\numexpr\csname djs@marks@#1\endcsname+#2\relax}}}
\RenewDocumentCommand{\marks}{o m}{%
  \protected@write\@auxout{}{\string\djsAddMarks{\number\value{enumi}}{#2}}%
  \IfValueTF{#1}%
    {\marginnote{\raggedright\textbf{[#2]}}[#1]}%
    {\marginnote{\raggedright\textbf{[#2]}}}%
}
\makeatother

\newcommand{\vspacerule}[2]{%
  \par\vspace*{#1}%
  \noindent\hrulefill\par
  \vspace*{#2}%
}

\definecolor{scarlet}{RGB}{205, 40, 75}
% Site colours for the contents-page title: the accent, the rule colour, a
% stronger rule colour for the path separator, and the muted text colour.
\definecolor{djsBurgundy}{RGB}{118, 59, 67}
\definecolor{djsHairline}{RGB}{216, 213, 205}
\definecolor{djsSlash}{RGB}{170, 166, 158}
\definecolor{djsMuted}{RGB}{96, 101, 106}

% Worked-solution box, adapted from the TMUA worked-solutions box. A white breakable box with a left scarlet rule and no surrounding
% frame or tint. A part that breaks fills the rest of its page, so the rule
% runs to the bottom margin; the unbroken or last part ends in a short
% horizontal foot rule so a solution visibly stops. Unlike the TMUA box there
% is no answer label. The box does not grow to the left, so the rule sits at
% the text edge; growing to the right keeps the usable line width close to the
% old tcolorbox Solution box.
% Spacing inside the box: paragraphs are 6pt apart, so source blank lines show
% (solutions put one after each display and none between a lead-in and its
% display). Every display gets 6pt above and 10pt below, whatever the length
% of its lead-in line; align rows are spread by an 8pt \jot; solution parts
% are 14pt apart and sub-parts 8pt; \left and \right space like ordinary
% brackets (mleftright); no single line is stranded at a page break; and
% neither the Solution label nor the line before a list ends a page. No
% skip can shrink, so full pages are not squeezed.
\newtcolorbox{djsSolution}{%
  enhanced, breakable, frame hidden, sharp corners, colback=white,
  height fixed for=first and middle, height fill,
  extras unbroken and last={natural height},
  borderline west={2.2pt}{0pt}{scarlet},
  left=11pt, right=0pt, top=5pt, bottom=13pt,
  before skip=13pt, after skip=4pt,
  grow to right by=10mm,
  overlay unbroken and last={%
    \fill[scarlet] (frame.south west)
      rectangle ([xshift=3.5cm, yshift=2.2pt]frame.south west);},
  before upper={{\color{scarlet}\bfseries Solution}\par\nobreak\vspace{4pt}%
    \setlength{\parskip}{6pt plus 2pt}%
    \setlength{\abovedisplayskip}{6pt plus 2pt}%
    \setlength{\abovedisplayshortskip}{6pt plus 2pt}%
    \setlength{\belowdisplayskip}{10pt plus 2pt}%
    \setlength{\belowdisplayshortskip}{10pt plus 2pt}%
    \setlength{\jot}{8pt}%
    \setlist{parsep=6pt plus 2pt, topsep=6pt plus 2pt, beginpenalty=10000}%
    \setlist[enumerate,1]{itemsep=14pt plus 2pt}%
    \setlist[enumerate,2]{itemsep=8pt plus 2pt}%
    \mleftright
    \clubpenalty=10000 \widowpenalty=10000 \displaywidowpenalty=10000\relax}}

% Optional smaller figure inside a solution box. \scalebox shrinks node text
% with the drawing. Not applied to existing figures.
\newcommand{\SolnFigureScale}{0.85}
\newsavebox{\djsSolnFigureBox}
\newenvironment{djsSolutionFigure}
  {\par\vspace{4pt}\begin{center}\begin{lrbox}{\djsSolnFigureBox}}
  {\end{lrbox}\scalebox{\SolnFigureScale}{\usebox{\djsSolnFigureBox}}%
   \end{center}\vspace{2pt}\par}

\newcommand{\cxl}[2]{%
  \tikz[baseline=(X.base)]{%
    \node[inner sep=0pt, outer sep=0pt, text=#1] (X) {$#2$};
    \draw[#1, line width=0.4pt] (X.south west)--(X.north east);
  }%
}

% Page-style hook run at the contents page and each question's opening page.
% A no-op for QBT packs; \djsSolnHeader points it at the fancy header.
\newcommand{\djsQuestionPageStyle}{}

% A contents entry: "Question N" in a box wide enough for two digits, so the
% mark totals line up, then the total in grey. Bookmarks keep "Question N".
\makeatletter
\newlength{\djsTocLabelWidth}
\newcommand{\djsTocMarks}[1]{%
  \@ifundefined{djs@marks@#1}{}{%
    {\footnotesize\itshape\color{black!45}%
     \csname djs@marks@#1\endcsname~mark%
     \ifnum\csname djs@marks@#1\endcsname=1 \else s\fi}}}
\newcommand{\djsTocEntry}[1]{%
  \texorpdfstring
    {\settowidth{\djsTocLabelWidth}{Question 00}%
     \makebox[\dimexpr\djsTocLabelWidth+0.9em\relax][l]{Question #1}%
     \djsTocMarks{#1}}%
    {Question #1}}
\makeatother

% The question count for the contents-page title is written to the .aux at
% the end of the document, so it also appears from the second pass.
\newcounter{djsQuestions}
\newcommand{\questionitem}{%
  \djsQuestionPageStyle
  \item\phantomsection
  \stepcounter{djsQuestions}%
  \addcontentsline{toc}{subsection}{\protect\djsTocEntry{\number\value{enumi}}}%
}
\makeatletter
\newcommand{\djsSetQuestionCount}[1]{\gdef\djs@questioncount{#1}}
\AtEndDocument{\immediate\write\@auxout
  {\string\djsSetQuestionCount{\number\value{djsQuestions}}}}
\makeatother

\renewcommand{\contentsname}{Questions}

% \djsPackHeader[<link content tag>]{<left header>}; the tag says which
% document a click on the header link came from.
\newcommand{\djsPackHeader}[2][questions]{%
  \pagestyle{fancy}%
  \fancyhead{}%
  \fancyfoot{}%
  \renewcommand{\headrulewidth}{0.9pt}%
  \lhead{#2}%
  \chead{}%
  \rhead{\href{https://danieljsmith.org/?utm_source=pdf&utm_campaign=further-maths&utm_content=#1}{danieljsmith.org}}%
}

% The header's topic is also the contents-page title, under which a label
% says which booklet this is.
\newcommand{\djsPackTitle}{}
\newcommand{\djsBookletLabel}{}
\newcommand{\djsQbtHeader}[1]{%
  \djsPackHeader{#1}%
  \renewcommand{\djsPackTitle}{#1}%
  \renewcommand{\djsBookletLabel}{Questions by topic}%
}

% The colophon at the foot of page one, independent of the page style: the
% contents-page title's hairline led by a short scarlet rule, then the date on
% the left and a link to the site on the right, all inside the bottom margin.
% Its sizes are fixed, so it matches the TMUA papers' colophon. The document
% supplies the date; recompiling never changes it.
\makeatletter
\newcommand{\djsLastUpdated}[1]{%
  \AtBeginDocument{\AddToHookNext{shipout/foreground}{\djs@colophon{#1}}}}
\newcommand{\djs@colophon}[1]{%
  \put(0.8in,-\dimexpr\paperheight-0.47in\relax){%
    \makebox[0pt][l]{\color{djsHairline}%
      \rule[0.9pt]{\dimexpr\paperwidth-1.8in\relax}{0.6pt}}%
    \makebox[0pt][l]{\color{scarlet}\rule{3.5cm}{2.2pt}}}%
  \put(0.8in,-\dimexpr\paperheight-0.32in\relax){\makebox[0pt][l]{%
    \fontsize{9}{11}\selectfont Last updated #1}}%
  \put(\dimexpr\paperwidth-1in\relax,-\dimexpr\paperheight-0.32in\relax){%
    \makebox[0pt][r]{\fontsize{7}{8.5}\selectfont\color{djsMuted}%
      Latest version at \href{https://resources.danieljsmith.org/}%
        {\color{black}resources.danieljsmith.org}}}}
\makeatother

% Solutions packs show the header only on the contents page and each
% question's opening page, so it never appears part-way through a solution.
% \thispagestyle marks the next page shipped out, which is the question's
% page because every solution question starts after a \newpage.
\newcommand{\djsSolnHeader}[1]{%
  \djsPackHeader[solutions]{\textit{Solutions} - #1}%
  \pagestyle{empty}%
  \renewcommand{\djsQuestionPageStyle}{\thispagestyle{fancy}}%
  \renewcommand{\djsPackTitle}{#1}%
  \renewcommand{\djsBookletLabel}{Worked solutions}%
}

% \djsSitePath{<strand>}{<topic>}, called once after the header, names where
% the pack sits on the site; the contents-page title prints it above the pack
% title, e.g. CORE PURE / FURTHER CALCULUS. Without it the line is left out.
\makeatletter
\newif\ifdjs@sitepath
\newcommand{\djs@strand}{}
\newcommand{\djs@section}{}
\newcommand{\djsSitePath}[2]{%
  \djs@sitepathtrue
  \renewcommand{\djs@strand}{#1}%
  \renewcommand{\djs@section}{#2}%
}
\makeatother

% Further Mechanics packs call \djsMechanicsPack in their preamble, so that
% \djsFrontMatter prints the default value of g above the contents. The box
% keeps its own space above it (before skip is the whole gap below the title,
% not an addition to it). Without the box, \djsGravityNote is just extra space
% before the contents heading, bringing the gap above it close to the gap
% between the box and the heading on Further Mechanics pages.
\newcommand{\djsGravityNote}{\vspace{14pt}}
\newcommand{\djsMechanicsPack}{%
  \renewcommand{\djsGravityNote}{%
    \begin{tcolorbox}[sharp corners, colback=white, colframe=black,
      boxrule=1pt, left=8pt, right=8pt, top=6pt, bottom=6pt,
      before skip=20pt]
    {\bfseries Instructions}\par\vspace{4pt}
    Unless a question states otherwise, take $g = 9.8\,\mathrm{m\,s^{-2}}$
    wherever a numerical value of $g$ is required. Solutions also give final
    answers for $g = 10\,\mathrm{m\,s^{-2}}$ where these differ.
    \end{tcolorbox}%
    \vspace{16pt}}%
}

% Contents-page title: the site path in spaced burgundy capitals, the pack
% title, a hairline led by a short scarlet rule (the solution box's foot rule),
% then the booklet label with the question and mark counts. The counts come
% from the .aux, so they appear from the second pass.
\makeatletter
\newsavebox{\djs@eyebrowbox}
\newcommand{\djs@caps}[1]{\textls[180]{\MakeUppercase{#1}}}
\newcommand{\djs@eyebrow}{%
  \color{djsBurgundy}\footnotesize\bfseries
  \djs@caps{\djs@strand}%
  {\color{djsSlash}\hspace{0.55em}/\hspace{0.55em}}%
  \djs@caps{\djs@section}}
\newcommand{\djs@dot}{\hspace{0.45em}\textperiodcentered\hspace{0.45em}}
\newcommand{\djs@counts}{%
  \@ifundefined{djs@questioncount}{}{%
    \djs@questioncount~question\ifnum\djs@questioncount=1 \else s\fi
    \ifnum\djs@totalmarks>0
      \djs@dot\the\djs@totalmarks~mark\ifnum\djs@totalmarks=1 \else s\fi
    \fi}}
\newcommand{\djsContentsTitle}{%
  \ifdjs@sitepath
    \sbox\djs@eyebrowbox{\djs@eyebrow}%
    \ifdim\wd\djs@eyebrowbox>\linewidth
      \@latex@warning@no@line{The \string\djsSitePath\space line is wider
        than the text, so it wraps. Shorten its topic}%
    \fi
    {\raggedright\djs@eyebrow\par}%
    \vspace{7pt}%
  \fi
  {\raggedright\fontsize{24.88}{28}\selectfont\djsPackTitle\par}%
  \vspace{9pt}%
  \noindent\makebox[0pt][l]{\color{djsHairline}\rule[0.9pt]{\linewidth}{0.6pt}}%
  {\color{scarlet}\rule{3.5cm}{2.2pt}}\par
  \vspace{5pt}%
  {\small\color{djsMuted}\djsBookletLabel\hfill\djs@counts\par}%
  \vspace{8pt}%
}
\makeatother

\newcommand{\djsFrontMatter}{%
  \djsQuestionPageStyle
  \vspace*{6pt}%
  \djsContentsTitle
  \djsGravityNote
  \tableofcontents
  \newpage
}


\djsSolnHeader{Motion Under a Variable Force}
\djsSitePath{Further Mechanics}{Kinematics \& Dynamics}
\djsMechanicsPack

\djsLastUpdated{27 October 2026}
\begin{document}
\djsFrontMatter
\begin{enumerate}[
  label=\textbf{\arabic*.},
  leftmargin=*,
  topsep=0pt,
  partopsep=0pt,
  parsep=0pt
]

% ---- Question 1 ----
\questionitem
% [Source: _QBT___Solns__Motion_Under_a_Variable_Force, Q1]

A particle $P$ of mass $4\,\mathrm{kg}$ moves along a straight horizontal line. At time $t\,\mathrm{s}$, its displacement in the positive direction from a fixed point $O$ is $x\,\mathrm{m}$ and its velocity is $v\,\mathrm{m\,s^{-1}}$. \\

The only horizontal force acting on $P$ has magnitude
\[
\dfrac{36}{(x+2)^3}\,\mathrm{N}
\] \\
and acts in the direction of decreasing $x$. Initially, $P$ is at $O$ and is moving in the positive direction with speed $\frac{3}{2}\,\mathrm{m\,s^{-1}}$. \\

\begin{enumerate}[label=\textbf{(\alph*)}]
  \item Find an expression for $v$ in terms of $x$. \marks{5} \\
  \item Hence find an expression for $x$ in terms of $t$. \marks{4} \\
\end{enumerate}

\vspace*{10pt}

\begin{djsSolution}
\begin{enumerate}[label=\textbf{(\alph*)}, start=1]
\item Take the direction of increasing \(x\) as positive. Using Newton's second law, with acceleration \(a=v\frac{\mathrm{d}v}{\mathrm{d}x}\):
\[
4v\frac{\mathrm{d}v}{\mathrm{d}x}=-\frac{36}{(x+2)^3}
\]

Separating the variables and integrating:
\begin{align*}
\int 4v\,\mathrm{d}v&=-36\int (x+2)^{-3}\,\mathrm{d}x\\
2v^2&=\frac{18}{(x+2)^2}+c
\end{align*}

Initially, \(x=0\) and \(v=\frac32\), so
\[
2\left(\frac32\right)^2=\frac{18}{2^2}+c
\quad\Longrightarrow\quad c=0
\]

Thus \(v^2=\frac{9}{(x+2)^2}\). The velocity is initially positive and cannot become zero at any finite \(x\geq 0\), so we take the positive root:
\[
v=\frac{3}{x+2}
\]

\item Using part (a) and \(v=\frac{\mathrm{d}x}{\mathrm{d}t}\), separate the variables and integrate:
\begin{align*}
\frac{\mathrm{d}x}{\mathrm{d}t}&=\frac{3}{x+2}\\
\int (x+2)\,\mathrm{d}x&=\int 3\,\mathrm{d}t\\
\frac{x^2}{2}+2x&=3t+d
\end{align*}

Since \(x=0\) when \(t=0\), \(d=0\). Completing the square gives
\begin{align*}
x^2+4x&=6t\\
(x+2)^2&=6t+4
\end{align*}

Taking the positive root since \(x+2>0\):
\[
x=\sqrt{6t+4}-2,\qquad t\geq 0
\]
\end{enumerate}
\end{djsSolution}

\newpage

% ---- Question 2 ----
\questionitem
% [Source: _QBT___Solns__Motion_Under_a_Variable_Force, Q2]

A lorry of mass $2000\,\mathrm{kg}$ travels along a straight horizontal road. Its engine delivers constant power of $80\,\mathrm{kW}$. When the lorry is moving at speed $v\,\mathrm{m\,s^{-1}}$, the resistance to its motion has magnitude $200v\,\mathrm{N}$. \\

At time $t$ seconds, the speed of the lorry is $v\,\mathrm{m\,s^{-1}}$. At $t=0$, the speed of the lorry is $8\,\mathrm{m\,s^{-1}}$. \\

\begin{enumerate}[label=\textbf{(\alph*)}]
  \item Show that
  \[
  10v\dfrac{\mathrm{d}v}{\mathrm{d}t}=400-v^2
  \]
  \marks[-30pt]{3}
  \item Verify that, while $8\leq v<20$,
  \[
  t=5\ln\left(\dfrac{336}{400-v^2}\right)
  \]
  \marks[-30pt]{4}
\end{enumerate}

\vspace*{10pt}

\begin{djsSolution}
\begin{enumerate}[label=\textbf{(\alph*)}, start=1]
  \item Take the direction of motion as positive. Using \(P=Fv\), the driving force is
  \[
  F=\frac{80\,000}{v}\,\mathrm{N}
  \]

  Using Newton's second law, with acceleration \(a=\frac{\mathrm{d}v}{\mathrm{d}t}\):
  \[
  2000\frac{\mathrm{d}v}{\mathrm{d}t}=\frac{80\,000}{v}-200v
  \]

  Multiplying by \(\frac{v}{200}\):
  \[
  10v\frac{\mathrm{d}v}{\mathrm{d}t}=400-v^2
  \]

  \item Separating the variables in the equation from part (a):
  \[
  \frac{10v}{400-v^2}\,\mathrm{d}v=\mathrm{d}t
  \]

  For \(8\leq v<20\), \(400-v^2>0\). Since the derivative of \(400-v^2\) is \(-2v\), integrating gives
  \begin{align*}
  \int 1\,\mathrm{d}t&=-5\int\frac{-2v}{400-v^2}\,\mathrm{d}v\\
  t&=-5\ln(400-v^2)+c
  \end{align*}

  Using \(v=8\) when \(t=0\):
  \begin{align*}
  0&=-5\ln(400-8^2)+c\\
  c&=5\ln336
  \end{align*}

  Substituting and combining the logarithms:
  \begin{align*}
  t&=5\ln336-5\ln(400-v^2)\\
  &=5\ln\left(\frac{336}{400-v^2}\right)
  \end{align*}
\end{enumerate}
\end{djsSolution}

\newpage

% ---- Question 3 ----
\questionitem
% [Source: _QBT___Solns__Motion_Under_a_Variable_Force, Q3]

A particle $P$ of mass $0.5\,\mathrm{kg}$ moves along the positive $x$-axis under the action of a single force. At time $t$ seconds, the particle is $x$ metres from the origin $O$, where $x>0$, and has speed $v\,\mathrm{m\,s^{-1}}$. The force acts in the direction of increasing $x$ and has magnitude
\[
\frac{k}{x^3}\text{ newtons}
\] \\
where $k$ is a positive constant. \\

When $x=2$, the speed of $P$ is $1\,\mathrm{m\,s^{-1}}$, and when $x=4$, its speed is $2\,\mathrm{m\,s^{-1}}$. \\

\begin{enumerate}[label=\textbf{(\alph*)}]
  \item Show that
  \[
  v^2=5-\frac{16}{x^2}
  \]
  \marks[-20pt]{5}
\end{enumerate}
The time taken for the speed of $P$ to increase from $1\,\mathrm{m\,s^{-1}}$ to $2\,\mathrm{m\,s^{-1}}$ is $T$ seconds. \\
\begin{enumerate}[label=\textbf{(\alph*)}, resume]
  \item Use algebraic integration to show that
  \[
  T=\frac{6}{5}
  \]
  \marks[-20pt]{4}
\end{enumerate}

\vspace*{10pt}

\begin{djsSolution}
\begin{enumerate}[label=\textbf{(\alph*)}, start=1]
\item Take the direction of increasing \(x\) as positive. Using Newton's second law, with acceleration \(a=v\frac{\mathrm{d}v}{\mathrm{d}x}\):
\[
\frac12 v\frac{\mathrm{d}v}{\mathrm{d}x}=\frac{k}{x^3}
\]

Separating the variables and integrating:
\begin{align*}
2v\,\mathrm{d}v&=4kx^{-3}\,\mathrm{d}x\\
\int 2v\,\mathrm{d}v&=\int 4kx^{-3}\,\mathrm{d}x\\
v^2&=-\frac{2k}{x^2}+c
\end{align*}

Substituting the two given conditions:
\begin{align*}
1&=-\frac{k}{2}+c\\
4&=-\frac{k}{8}+c
\end{align*}

Subtracting the first equation from the second:
\begin{align*}
3&=\frac{3k}{8}\\
k&=8\\
c&=1+\frac{8}{2}=5
\end{align*}

Hence
\[
v^2=5-\frac{16}{x^2}
\]

\item Since the particle moves in the direction of increasing \(x\), use the positive square root from part (a). As \(x>0\), the velocity is
\[
\frac{\mathrm{d}x}{\mathrm{d}t}=v=\sqrt{5-\frac{16}{x^2}}=\frac{\sqrt{5x^2-16}}{x}
\]

Separating the variables and integrating:
\begin{align*}
\mathrm{d}t&=\frac{x}{\sqrt{5x^2-16}}\,\mathrm{d}x\\
\int 1\,\mathrm{d}t&=\frac{1}{10}\int 10x(5x^2-16)^{-1/2}\,\mathrm{d}x\\
t&=\frac15\sqrt{5x^2-16}+d
\end{align*}

Let \(t_1\) and \(t_2\) be the times when \(x=2\) and \(x=4\), respectively. Substituting these values:
\begin{align*}
t_1&=\frac15\sqrt{20-16}+d=\frac25+d\\
t_2&=\frac15\sqrt{80-16}+d=\frac85+d\\
T&=t_2-t_1=\frac85-\frac25=\frac65
\end{align*}
\end{enumerate}
\end{djsSolution}

\newpage

% ---- Question 4 ----
\questionitem
% [Source: _QBT___Solns__Motion_Under_a_Variable_Force, Q4]

A particle $P$ of mass $3\,\mathrm{kg}$ moves along the positive $x$-axis. At time $t=0$, it passes through the origin with speed $4\,\mathrm{m\,s^{-1}}$. \\

While $P$ is moving, the only force acting on it is a resistive force of magnitude
\[
\frac{6}{v}\,\mathrm{N}
\] \\
where $v\,\mathrm{m\,s^{-1}}$ is its speed. The force acts in the negative $x$-direction. \\

\begin{enumerate}[label=\textbf{(\alph*)}]
  \item Find an expression for $v$ in terms of the displacement $x$ of $P$ from the origin. \marks{5} \\
  \item Find the distance travelled by $P$ while its speed decreases from $4\,\mathrm{m\,s^{-1}}$ to $2\,\mathrm{m\,s^{-1}}$. \marks{2} \\
\end{enumerate}
A particle $Q$, also of mass $3\,\mathrm{kg}$, starts from the origin with speed $4\,\mathrm{m\,s^{-1}}$. The only force acting on $Q$ is a constant resistive force of magnitude $\frac{3}{2}\,\mathrm{N}$. \\
\begin{enumerate}[label=\textbf{(\alph*)}, resume]
  \item By comparing the resistive forces acting on $P$ and $Q$, explain why $P$ must come to rest at a time $t<8\,\mathrm{s}$ and at a displacement $x<16\,\mathrm{m}$. \marks{3} \\
  \item Sketch the velocity--time graph for $P$ from $t=0$ until it first comes to rest. Your sketch should indicate the concavity of the graph, but no numerical values are required. \marks{1} \\
  \item Determine the maximum displacement of $P$ from the origin. \marks{2} \\
\end{enumerate}

\vspace*{10pt}

\begin{djsSolution}
\begin{enumerate}[label=\textbf{(\alph*)}, start=1]
\item Take the direction of increasing \(x\) as positive. Using Newton's second law, with acceleration \(a=v\frac{\mathrm{d}v}{\mathrm{d}x}\):
\begin{align*}
3v\frac{\mathrm{d}v}{\mathrm{d}x}&=-\frac{6}{v}\\
v^2\frac{\mathrm{d}v}{\mathrm{d}x}&=-2
\end{align*}

Separating the variables and integrating:
\begin{align*}
\int v^2\,\mathrm{d}v&=-2\int 1\,\mathrm{d}x\\
\frac{v^3}{3}&=-2x+c
\end{align*}

When \(x=0\), \(v=4\), so \(c=\frac{64}{3}\). Hence
\begin{align*}
v^3&=64-6x\\
v&=\sqrt[3]{64-6x}
\end{align*}

\item Substituting \(v=2\) into the result of part (a):
\[
8=64-6x
\quad\Longrightarrow\quad
x=\frac{28}{3}\,\mathrm{m}
\]

Since \(P\) moves in the positive direction throughout, the distance travelled is \(\frac{28}{3}\,\mathrm{m}\).

\item For \(Q\), Newton's second law in the same positive direction gives
\[
3a_Q=-\frac32
\quad\Longrightarrow\quad
a_Q=-\frac12\,\mathrm{m\,s^{-2}}
\]

Using constant-acceleration formulae, its stopping time and displacement are
\begin{align*}
0&=4-\frac12t
\quad\Longrightarrow\quad t=8\,\mathrm{s}\\
x_Q&=4(8)+\frac12\left(-\frac12\right)(8)^2=16\,\mathrm{m}
\end{align*}

The initial resistive forces are equal. As \(P\) slows down, however, its resistive force increases:
\[
0<v<4
\quad\Longrightarrow\quad
\frac6v>\frac32
\]

Thus, at every speed below \(4\,\mathrm{m\,s^{-1}}\), \(P\) has a greater deceleration than \(Q\). Each further decrease in speed therefore takes less time and covers less distance for \(P\), so it comes to rest with
\[
t<8\,\mathrm{s},\qquad x<16\,\mathrm{m}
\]

\item The sketch should look like:\par\nopagebreak

\begin{tikzpicture}[x=0.12\linewidth,y=2cm]
\draw[->] (0,0) -- (6,0) node[right] {\(t\)};
\draw[->] (0,0) -- (0,4.1) node[above] {\(v\)};
\draw[thick] (0,3.6) .. controls (3.333,2.4) and (5,1.2) .. (5,0);
\end{tikzpicture}

As the speed falls, the resistive force increases, making the acceleration more negative. The graph is therefore decreasing and concave.

\item The displacement increases while \(P\) is moving, so its maximum occurs as \(P\) comes to rest. Taking \(v\to0\) in part (a):
\[
0=64-6x_{\max}
\quad\Longrightarrow\quad
x_{\max}=\frac{32}{3}\,\mathrm{m}
\]
\end{enumerate}
\end{djsSolution}

\newpage

% ---- Question 5 ----
\questionitem
% [Source: _QBT___Solns__Motion_Under_a_Variable_Force, Q5]

A particle $B$ of mass $2\,\mathrm{kg}$ moves along the positive $x$-axis. When $B$ passes through the origin, its speed is $u\,\mathrm{m\,s^{-1}}$. \\

The only horizontal force acting on $B$ is a resistance whose magnitude is directly proportional to the cube of the speed of $B$. When the speed of $B$ is $3\,\mathrm{m\,s^{-1}}$, the magnitude of the resistance is $54\,\mathrm{N}$. \\

\begin{enumerate}[label=\textbf{(\alph*)}]
  \item Show that, while $B$ is moving in the positive $x$-direction, its motion satisfies
  \[
  \frac{1}{v^2}\frac{\mathrm{d}v}{\mathrm{d}x}=-1
  \]
  \marks[-30pt]{3}
  \item
  \begin{enumerate}[label=(\roman*)]
    \item Solve the differential equation in part (a) to obtain $v$ in terms of $x$ and $u$. \marks{4} \\
    \item By considering the behaviour of $v$ as $x\to\infty$, state one respect in which this model of the motion is unrealistic. \marks{1} \\
  \end{enumerate}
\end{enumerate}
The particle subsequently reaches a point $A$ with coordinate $X$, where its speed is $V\,\mathrm{m\,s^{-1}}$. The work done by the resistance as $B$ moves from the origin to $A$ is $W\,\mathrm{J}$. \\
\begin{enumerate}[label=\textbf{(\alph*)}, resume]
  \item
  \begin{enumerate}[label=(\roman*)]
    \item Using $W=\int F\,\mathrm{d}x$, where $F$ is the signed force in the $x$-direction, find an expression for $W$ in terms of $X$ and $u$. \marks{3} \\
    \item Explain the physical significance of the sign of $W$. \marks{1} \\
    \item Rewrite your answer to part (c)(i) in terms of $V$ and $u$, and hence show how $W$ is related to the change in the kinetic energy of $B$. \marks{2} \\
  \end{enumerate}
\end{enumerate}

\vspace*{10pt}

\begin{djsSolution}
\begin{enumerate}[label=\textbf{(\alph*)}, start=1]
  \item Take the direction of increasing \(x\) as positive. Let the magnitude of the resistance be \(kv^3\). The given force and speed give
  \[
  54=k(3^3)\quad\Longrightarrow\quad k=2
  \]

  Using Newton's second law, with acceleration \(a=v\frac{\mathrm{d}v}{\mathrm{d}x}\):
  \[
  2v\frac{\mathrm{d}v}{\mathrm{d}x}=-kv^3=-2v^3
  \]

  While the particle is moving in the positive direction, \(v>0\), so dividing by \(2v^3\) gives
  \[
  \frac{1}{v^2}\frac{\mathrm{d}v}{\mathrm{d}x}=-1
  \]

  \item
  \begin{enumerate}[label=(\roman*), start=1]
    \item Separating the variables and integrating:
    \begin{align*}
    \int v^{-2}\,\mathrm{d}v&=-\int 1\,\mathrm{d}x\\
    -\frac{1}{v}&=-x+c
    \end{align*}

    At the origin, \(x=0\) and \(v=u\), so \(c=-\frac{1}{u}\). Hence
    \begin{align*}
    -\frac{1}{v}&=-x-\frac{1}{u}\\
    \frac{1}{v}&=x+\frac{1}{u}=\frac{1+ux}{u}\\
    v&=\frac{u}{1+ux}
    \end{align*}

    \item As \(x\to\infty\), \(v\to0\), but \(v>0\) for every finite \(x\). Thus the model unrealistically predicts that the particle travels an unlimited distance without coming to rest.
  \end{enumerate}

  \item
  \begin{enumerate}[label=(\roman*), start=1]
    \item The signed force is \(F=-2v^3\). Using the speed from part (b)(i):
    \begin{align*}
    W&=\int_0^X F\,\mathrm{d}x\\
    &=-2u^3\int_0^X(1+ux)^{-3}\,\mathrm{d}x\\
    &=-2u^3\left[\frac{(1+ux)^{-2}}{-2u}\right]_0^X\\
    &=u^2\left[(1+uX)^{-2}-1\right]
    \end{align*}

    \item \(W<0\): the resistance acts opposite to the displacement, so it does negative work and reduces the particle's kinetic energy.

    \item From part (b)(i), \(V=\frac{u}{1+uX}\), so the expression in part (c)(i) becomes
    \[
    W=\frac{u^2}{(1+uX)^2}-u^2=V^2-u^2
    \]

    Since the mass is \(2\,\mathrm{kg}\), the change in kinetic energy is
    \[
    \frac12(2)V^2-\frac12(2)u^2=V^2-u^2=W
    \]
  \end{enumerate}
\end{enumerate}
\end{djsSolution}

\newpage

% ---- Question 6 ----
\questionitem
% [Source: _QBT___Solns__Motion_Under_a_Variable_Force, Q6]

A small sphere $P$ of mass $5\,\mathrm{kg}$ is released into a liquid and moves vertically downwards. Take the downward direction as positive. In a model of the motion, the forces acting on $P$ are: \\

\begin{itemize}
  \item its weight;
  \item a constant upthrust of magnitude $14\,\mathrm{N}$;
  \item a resistance to motion whose magnitude is proportional to the speed of $P$.
\end{itemize}

When the speed of $P$ is $4\,\mathrm{m\,s^{-1}}$, the resistance has magnitude $10\,\mathrm{N}$. At time $t=0$, the sphere is at a point $O$ and is moving downwards with speed $u\,\mathrm{m\,s^{-1}}$, where $u>0$. At time $t$ seconds, its downward velocity is $v\,\mathrm{m\,s^{-1}}$. Take $g=9.8\,\mathrm{m\,s^{-2}}$. \\

\begin{enumerate}[label=\textbf{(\alph*)}]
  \item Show that, according to the model,
  \[
  \dfrac{\mathrm{d}v}{\mathrm{d}t}=7-\dfrac{1}{2}v
  \]
  \marks[-30pt]{4}
  \item
  \begin{enumerate}[label=(\roman*)]
    \item Solve the differential equation in part (a) to show that
    \[
    v=14+(u-14)e^{-t/2}
    \]
    \marks[-20pt]{5}
    \item Describe briefly how the speed of $P$ varies over time in each of the following cases:
    \begin{itemize}
      \item $0<u<14$;
      \item $u=14$;
      \item $u>14$.
    \end{itemize}
    In each case, state the limiting speed of $P$ as $t\to\infty$. \marks{3} \\
  \end{enumerate}
  \item Given that $u=9$, find the exact downward displacement of $P$ from $O$ when $t=6$. \marks{4} \\
\end{enumerate}

\vspace*{10pt}

\begin{djsSolution}
\begin{enumerate}[label=\textbf{(\alph*)}, start=1]
\item Write the resistance as \(kv\). Since \(4k=10\), \(k=\frac52\).

Using Newton's second law, with acceleration \(a=\frac{\mathrm{d}v}{\mathrm{d}t}\):
\begin{align*}
5\dfrac{\mathrm{d}v}{\mathrm{d}t}&=5(9.8)-14-\dfrac52v\\
\dfrac{\mathrm{d}v}{\mathrm{d}t}&=7-\dfrac12v
\end{align*}

\item
\begin{enumerate}[label=(\roman*), start=1]
\item For \(u\ne14\), separate the variables in the equation from part (a):
\begin{align*}
\dfrac{\mathrm{d}v}{\mathrm{d}t}&=-\dfrac12(v-14)\\
\int\dfrac{1}{v-14}\,\mathrm{d}v&=-\dfrac12\int1\,\mathrm{d}t\\
\ln|v-14|&=-\dfrac t2+c
\end{align*}

Using \(v=u\) when \(t=0\) gives \(c=\ln|u-14|\). Hence
\begin{align*}
\ln|v-14|-\ln|u-14|&=-\dfrac t2\\
|v-14|&=|u-14|e^{-t/2}
\end{align*}

The right-hand side is non-zero for finite \(t\), so \(v-14\) cannot change sign. It therefore has the same sign as its initial value \(u-14\), giving
\[
v-14=(u-14)e^{-t/2}
\]

If \(u=14\), the differential equation gives the constant solution \(v=14\). Thus, for every \(u>0\), the solution is
\[
v=14+(u-14)e^{-t/2}
\]

\item Since \(e^{-t/2}\) decreases from \(1\) towards \(0\):
\begin{itemize}
\item If \(0<u<14\), the speed increases from \(u\) towards \(14\,\mathrm{m\,s^{-1}}\).
\item If \(u=14\), the speed remains constant at \(14\,\mathrm{m\,s^{-1}}\).
\item If \(u>14\), the speed decreases from \(u\) towards \(14\,\mathrm{m\,s^{-1}}\).
\end{itemize}
In all three cases, the limiting speed is \(14\,\mathrm{m\,s^{-1}}\).
\end{enumerate}

\item Let \(x\) be the downward displacement from \(O\), so \(x=0\) when \(t=0\). Using \(u=9\) in part (b)(i):
\[
\dfrac{\mathrm{d}x}{\mathrm{d}t}=v=14-5e^{-t/2}
\]

Integrating with respect to \(t\):
\begin{align*}
x&=14t+10e^{-t/2}+d\\
0&=10+d\\
d&=-10
\end{align*}

Therefore, at \(t=6\), the downward displacement is
\[
x=14(6)+10e^{-3}-10=\left(74+10e^{-3}\right)\,\mathrm{m}
\]
\end{enumerate}
\end{djsSolution}

\newpage

% ---- Question 7 ----
\questionitem
% [Source: _QBT___Solns__Motion_Under_a_Variable_Force, Q7]

A vehicle of mass $900\,\mathrm{kg}$ moves along a straight horizontal road. Its engine provides a constant driving force of $576\,\mathrm{N}$. When the vehicle has speed $v\,\mathrm{m\,s^{-1}}$, the resistance to its motion has magnitude $(v^2+14v)\,\mathrm{N}$. \\

The time taken for the speed of the vehicle to increase from $3\,\mathrm{m\,s^{-1}}$ to $13\,\mathrm{m\,s^{-1}}$ is $T$ seconds. \\

\begin{enumerate}[label=\textbf{(\alph*)}]
  \item Show that
  \[
  T=18\ln\left(\frac{27}{7}\right)
  \]
  \marks[-30pt]{6}
\end{enumerate}
When $t=0$, the speed of the vehicle is $8\,\mathrm{m\,s^{-1}}$. \\
\begin{enumerate}[label=\textbf{(\alph*)}, resume]
  \item Show that
  \[
  \frac{v+32}{18-v}=pe^{qt}
  \]
  where the rational numbers $p$ and $q$ are to be found. \marks{2} \\
  \item Hence explain why $v<18$ for all $t\geq 0$. \marks{1} \\
\end{enumerate}

\vspace*{10pt}

\begin{djsSolution}
\begin{enumerate}[label=\textbf{(\alph*)}, start=1]
\item Take the direction of motion as positive. Using Newton's second law, with acceleration \(a=\frac{\mathrm{d}v}{\mathrm{d}t}\):
\begin{align*}
900\frac{\mathrm{d}v}{\mathrm{d}t}
&=576-(v^2+14v)\\
&=(18-v)(v+32)
\end{align*}

Separating the variables:
\[
\frac{900}{(18-v)(v+32)}\,\mathrm{d}v=\mathrm{d}t
\]

For the partial fractions, write
\[
\frac{1}{(18-v)(v+32)}=\frac{A}{18-v}+\frac{B}{v+32}
\]

Multiplying by \((18-v)(v+32)\) and comparing coefficients gives
\begin{align*}
1&=A(v+32)+B(18-v)\\
&=(A-B)v+32A+18B\\
A-B&=0,\qquad 32A+18B=1\\
A&=B=\frac{1}{50}
\end{align*}

Both \(18-v\) and \(v+32\) are positive while \(0\leq v<18\), which part (c) shows holds throughout the motion. Integrating both sides:
\begin{align*}
18\int\left(\frac{1}{18-v}+\frac{1}{v+32}\right)\,\mathrm{d}v
&=\int 1\,\mathrm{d}t\\
18\left(-\ln(18-v)+\ln(v+32)\right)&=t+c
\end{align*}

Combining the logarithms:
\[
18\ln\left(\frac{v+32}{18-v}\right)=t+c
\]

Let \(t_3\) and \(t_{13}\) be the times at which the speeds are \(3\) and \(13\,\mathrm{m\,s^{-1}}\), respectively. Substituting these speeds:
\begin{align*}
t_3+c&=18\ln\left(\frac{35}{15}\right)=18\ln\left(\frac73\right)\\
t_{13}+c&=18\ln\left(\frac{45}{5}\right)=18\ln9
\end{align*}

Subtracting eliminates \(c\):
\begin{align*}
T=t_{13}-t_3
&=18\left(\ln9-\ln\left(\frac73\right)\right)\\
&=18\ln\left(\frac{27}{7}\right)
\end{align*}

\item Using the integrated relation from part (a), the condition \(v=8\) when \(t=0\) gives
\[
c=18\ln\left(\frac{8+32}{18-8}\right)=18\ln4
\]

Thus, rearranging and exponentiating:
\begin{align*}
\ln\left(\frac{v+32}{18-v}\right)&=\frac{t}{18}+\ln4\\
\frac{v+32}{18-v}&=4e^{t/18}
\end{align*}

This is the required form \(pe^{qt}\), with
\[
p=4,\qquad q=\frac{1}{18}
\]

\item From part (b), \(v+32=4e^{t/18}(18-v)\). Since \(v+32=50-(18-v)\), this rearranges to
\[
18-v=\frac{50}{1+4e^{t/18}}>0
\]

The denominator is positive for every \(t\geq0\), so \(v<18\) for all \(t\geq0\).
\end{enumerate}
\end{djsSolution}

\newpage

% ---- Question 8 ----
\questionitem
% [Source: _QBT___Solns__Motion_Under_a_Variable_Force, Q8]

A particle $P$ of mass $2\,\mathrm{kg}$ moves along the positive $x$-axis in the direction of increasing $x$. At time $t$ seconds, its displacement from the origin is $x$ metres and its speed is $v\,\mathrm{m\,s^{-1}}$. \\

While $P$ is moving, the resultant force acting on it is directed towards the origin and has magnitude $kv(v+2)\,\mathrm{N}$, where $k$ is a positive constant and $v$ denotes the numerical value of the speed in metres per second. \\

When $x=1$, the speed of $P$ is $7\,\mathrm{m\,s^{-1}}$, and when $x=3$, its speed is $1\,\mathrm{m\,s^{-1}}$. \\

\begin{enumerate}[label=\textbf{(\alph*)}]
  \item Find $k$ and show that
  \[
  v=9\times 3^{(1-x)/2}-2
  \]
  \marks[-20pt]{6}
\end{enumerate}
The time taken for the speed of $P$ to decrease from $7\,\mathrm{m\,s^{-1}}$ to $1\,\mathrm{m\,s^{-1}}$ is $T$ seconds. \\
\begin{enumerate}[label=\textbf{(\alph*)}, resume]
  \item Show that
  \[
  T=\frac{\ln(7/3)}{\ln 3}
  \]
  \marks[-30pt]{5}
\end{enumerate}

\vspace*{10pt}

\begin{djsSolution}
\begin{enumerate}[label=\textbf{(\alph*)}, start=1]
\item Take the direction of increasing \(x\) as positive. The resultant force acts in the negative direction. Using Newton's second law, with acceleration \(a=v\frac{\mathrm{d}v}{\mathrm{d}x}\):
\[
2v\frac{\mathrm{d}v}{\mathrm{d}x}=-kv(v+2)
\]

Since \(v>0\) while the particle is moving, \(v+2>0\). Dividing by \(v\) and separating the variables:
\begin{align*}
\frac{1}{v+2}\,\mathrm{d}v&=-\frac{k}{2}\,\mathrm{d}x\\
\int\frac{1}{v+2}\,\mathrm{d}v&=-\frac{k}{2}\int 1\,\mathrm{d}x\\
\ln(v+2)&=-\frac{kx}{2}+c
\end{align*}

Substituting the two given conditions:
\begin{align*}
\ln 9&=-\frac{k}{2}+c\\
\ln 3&=-\frac{3k}{2}+c
\end{align*}

Subtracting the first equation from the second:
\begin{align*}
\ln 3-\ln 9&=-k\\
k&=\ln 3
\end{align*}

Subtracting the relation at \(x=1\), \(v=7\) from the general relation now gives
\begin{align*}
\ln(v+2)-\ln 9&=-\frac{k}{2}(x-1)\\
\ln\left(\frac{v+2}{9}\right)&=\frac{1-x}{2}\ln 3\\
\frac{v+2}{9}&=3^{(1-x)/2}\\
v&=9\times 3^{(1-x)/2}-2
\end{align*}

\item Using Newton's second law, with acceleration \(a=\frac{\mathrm{d}v}{\mathrm{d}t}\):
\[
2\frac{\mathrm{d}v}{\mathrm{d}t}=-kv(v+2)
\]

Both \(v\) and \(v+2\) are positive. Separating the variables, and using partial fractions:
\begin{align*}
\frac{2}{v(v+2)}\,\mathrm{d}v&=-k\,\mathrm{d}t\\
\frac{2}{v(v+2)}&=\frac{1}{v}-\frac{1}{v+2}\\
\int\left(\frac{1}{v}-\frac{1}{v+2}\right)\,\mathrm{d}v&=-k\int 1\,\mathrm{d}t\\
\ln v-\ln(v+2)&=-kt+d
\end{align*}

Combining the logarithms:
\[
\ln\left(\frac{v}{v+2}\right)=-kt+d
\]

Let \(t_0\) be the time when \(v=7\). Then \(v=1\) at time \(t_0+T\), giving
\begin{align*}
\ln\left(\frac{7}{9}\right)&=-kt_0+d\\
\ln\left(\frac{1}{3}\right)&=-k(t_0+T)+d
\end{align*}

Subtracting eliminates \(d\) and \(t_0\). Using \(k=\ln 3\) from part (a):
\begin{align*}
\ln\left(\frac{1}{3}\right)-\ln\left(\frac{7}{9}\right)&=-kT\\
\ln\left(\frac{3}{7}\right)&=-T\ln 3\\
T&=\frac{\ln(7/3)}{\ln 3}
\end{align*}
\end{enumerate}
\end{djsSolution}

\newpage

% ---- Question 9 ----
\questionitem
% [Source: _QBT___Solns__Motion_Under_a_Variable_Force, Q9]

A particle $P$ of mass $4\,\mathrm{kg}$ is released from rest and moves vertically downwards through a fluid. While $P$ is moving with speed $v\,\mathrm{m\,s^{-1}}$, it experiences a resistance of magnitude $6v\,\mathrm{N}$ acting vertically upwards. After $P$ has descended a distance $x\,\mathrm{m}$, its speed is $v\,\mathrm{m\,s^{-1}}$. Take $g=9.8\,\mathrm{m\,s^{-2}}$. \\

\begin{enumerate}[label=\textbf{(\alph*)}]
  \item Show that
  \[
  \left(\frac{10v}{98-15v}\right)\frac{\mathrm{d}v}{\mathrm{d}x}=1
  \]
  \marks[-30pt]{2}
  \item Verify that $x$ and $v$ satisfy
  \[
  45x=196\ln\left(\frac{98}{98-15v}\right)-30v
  \]
  \marks[-30pt]{5}
  \item Determine the work done against the resistance from the instant $P$ is released until its acceleration has magnitude $2.45\,\mathrm{m\,s^{-2}}$. Give your answer to $3$ significant figures. \marks{5} \\
\end{enumerate}

\vspace*{10pt}

\begin{djsSolution}
\begin{enumerate}[label=\textbf{(\alph*)}, start=1]
\item Take vertically downwards as positive. Using Newton's second law, with acceleration \(a=v\frac{\mathrm{d}v}{\mathrm{d}x}\):
\begin{align*}
4v\frac{\mathrm{d}v}{\mathrm{d}x}&=4(9.8)-6v\\
10v\frac{\mathrm{d}v}{\mathrm{d}x}&=98-15v\\
\left(\frac{10v}{98-15v}\right)\frac{\mathrm{d}v}{\mathrm{d}x}&=1
\end{align*}

\item Separating the variables in the equation from part (a):
\[
\int\frac{10v}{98-15v}\,\mathrm{d}v=\int 1\,\mathrm{d}x
\]

The terminal speed is \(\frac{98}{15}\,\mathrm{m\,s^{-1}}\). Since the particle accelerates towards this speed from rest, \(98-15v>0\). To integrate, write
\[
\frac{10v}{98-15v}=-\frac23+\frac{196}{3(98-15v)}
\]

Hence
\[
x=-\frac{2v}{3}-\frac{196}{45}\ln(98-15v)+c
\]

Initially, \(x=0\) and \(v=0\), so
\[
0=-\frac{196}{45}\ln98+c
\quad\Longrightarrow\quad
c=\frac{196}{45}\ln98
\]

Substituting for \(c\) and combining the logarithms:
\begin{align*}
45x&=-30v-196\ln(98-15v)+196\ln98\\
&=196\ln\left(\frac{98}{98-15v}\right)-30v
\end{align*}

\item The acceleration is downwards throughout the motion. Newton's second law gives
\[
a=9.8-1.5v
\]

Setting its magnitude to \(2.45\,\mathrm{m\,s^{-2}}\):
\[
2.45=9.8-1.5v
\quad\Longrightarrow\quad
v=4.9\,\mathrm{m\,s^{-1}}
\]

Using part (b), the distance descended is
\begin{align*}
45x&=196\ln\left(\frac{98}{98-15(4.9)}\right)-30(4.9)\\
x&=\frac{196\ln4-147}{45}\,\mathrm{m}
\end{align*}

Let \(W\) be the positive work done against resistance. The work done by resistance is \(-W\). The work--energy principle, with work positive when a force acts in the direction of motion, is
\[
KE_0+GPE_0+\sum W_{\text{other}}
=KE_1+GPE_1
\]

Take the final position as the zero level for gravitational potential energy. Substituting:
\begin{align*}
0+4gx-W&=\frac12(4)v^2+0\\
W&=39.2x-2v^2\\
&=39.2\left(\frac{196\ln4-147}{45}\right)-2(4.9)^2\\
&\approx60.6\,\mathrm{J}\qquad\text{(3 s.f.)}
\end{align*}
\end{enumerate}
\end{djsSolution}

\newpage

% ---- Question 10 ----
\questionitem
% [Source: _QBT___Solns__Motion_Under_a_Variable_Force, Q10]

A particle $P$, of mass $m$, moves on a rough straight track inclined at an angle $\alpha$ to the horizontal. A fixed point $O$ lies on the track, and $x$ denotes the distance of $P$ from $O$, measured up the track. \\

In addition to its weight, $P$ is acted on by a force directed down the track and having magnitude
\[
\frac{kmg}{x+a}
\] \\
where $a$ and $k$ are positive constants with dimensions of length. The coefficient of friction between $P$ and the track is $\mu$. \\

The particle is projected up the track and first comes to rest at the point where $x=2a$. \\

\begin{enumerate}[label=\textbf{(\alph*)}]
  \item Show that, while $P$ is moving up the track with speed $v$,
  \[
  \frac{\mathrm{d}}{\mathrm{d}x}(v^2)+\frac{2gk}{x+a}+2g(\sin\alpha+\mu\cos\alpha)=0
  \]
  \marks[-30pt]{4}
  \item Verify that
  \[
  v^2=2gk\ln\left(\frac{3a}{x+a}\right)+2g(\sin\alpha+\mu\cos\alpha)(2a-x)
  \]
  \marks[-30pt]{4}
  \item Find, in terms of $a$, $k$ and $\alpha$, the range of values of $\mu$ for which $P$ remains at rest when $x=2a$. \marks{2} \\
\end{enumerate}

\vspace*{10pt}

\begin{djsSolution}
\begin{enumerate}[label=\textbf{(\alph*)}, start=1]
\item Take up the track as positive. Let the normal reaction be \(R\). Resolving perpendicular to the track gives
\[
R=mg\cos\alpha
\]

While \(P\) moves up the track, friction acts down the track with magnitude \(\mu R\). Using Newton's second law, with acceleration \(v\frac{\mathrm{d}v}{\mathrm{d}x}\):
\begin{align*}
mv\frac{\mathrm{d}v}{\mathrm{d}x}
&=-mg\sin\alpha-\mu R-\frac{kmg}{x+a}\\
2v\frac{\mathrm{d}v}{\mathrm{d}x}
&=-2g(\sin\alpha+\mu\cos\alpha)-\frac{2gk}{x+a}
\end{align*}

Since \(\frac{\mathrm{d}}{\mathrm{d}x}(v^2)=2v\frac{\mathrm{d}v}{\mathrm{d}x}\), this becomes
\[
\frac{\mathrm{d}}{\mathrm{d}x}(v^2)+\frac{2gk}{x+a}+2g(\sin\alpha+\mu\cos\alpha)=0
\]

\item Since \(x\) is a distance and \(a>0\), \(x+a>0\). Integrating the equation in part (a) with respect to \(x\):
\[
v^2=-2gk\ln(x+a)-2g(\sin\alpha+\mu\cos\alpha)x+c
\]

Substituting \(v=0\) when \(x=2a\) gives
\begin{align*}
0&=-2gk\ln(3a)-4ga(\sin\alpha+\mu\cos\alpha)+c\\
c&=2gk\ln(3a)+4ga(\sin\alpha+\mu\cos\alpha)
\end{align*}

Hence
\begin{align*}
v^2&=2gk\bigl(\ln(3a)-\ln(x+a)\bigr)
+2g(\sin\alpha+\mu\cos\alpha)(2a-x)\\
&=2gk\ln\left(\frac{3a}{x+a}\right)
+2g(\sin\alpha+\mu\cos\alpha)(2a-x)
\end{align*}

\item At \(x=2a\), let the frictional force acting up the track be \(F\). For equilibrium, Newton's second law gives
\[
F-mg\sin\alpha-\frac{kmg}{3a}=0
\]

The required friction must not exceed its limiting value \(\mu R=\mu mg\cos\alpha\), so
\begin{align*}
mg\sin\alpha+\frac{kmg}{3a}&\leq\mu mg\cos\alpha\\
\mu&\geq\tan\alpha+\frac{k}{3a\cos\alpha}
\end{align*}
\end{enumerate}
\end{djsSolution}

\newpage

% ---- Question 11 ----
\questionitem
% [Source: _QBT___Solns__Motion_Under_a_Variable_Force, Q11]

A small instrument package of mass $m\,\mathrm{kg}$ enters a deep liquid moving vertically downwards at $7\,\mathrm{m\,s^{-1}}$. While the package is moving downwards, it experiences a resistive force of magnitude $\frac{m}{20}v^2\,\mathrm{N}$, where $v\,\mathrm{m\,s^{-1}}$ is its speed. \\

At time $t\,\mathrm{s}$ after entering the liquid, the package has descended a distance $x\,\mathrm{m}$. \\

\begin{enumerate}[label=\textbf{(\alph*)}]
  \item Find an expression for $v$ in terms of $t$. \marks{5} \\
  \item Find an expression for $x$ in terms of $t$. \marks{4} \\
  \item Find the exact distance descended by the package when its speed first reaches $\frac{21}{2}\,\mathrm{m\,s^{-1}}$. \marks{4} \\
\end{enumerate}

\vspace*{10pt}

\begin{djsSolution}
\begin{enumerate}[label=\textbf{(\alph*)}, start=1]
\item Take vertically downwards as positive. Using Newton's second law, with acceleration \(a=\frac{\mathrm{d}v}{\mathrm{d}t}\):
\begin{align*}
m\frac{\mathrm{d}v}{\mathrm{d}t} &= mg-\frac{m}{20}v^2 \\
20\frac{\mathrm{d}v}{\mathrm{d}t} &= 196-v^2
\end{align*}

The terminal speed is \(14\,\mathrm{m\,s^{-1}}\). Since the initial speed is \(7\,\mathrm{m\,s^{-1}}\), the speed increases towards \(14\,\mathrm{m\,s^{-1}}\) and remains below it, so \(14+v\) and \(14-v\) are both positive.

Separating the variables and using partial fractions:
\begin{align*}
\int\frac{1}{196-v^2}\,\mathrm{d}v &= \frac{1}{20}\int 1\,\mathrm{d}t \\
\frac{1}{28}\int\left(\frac{1}{14-v}+\frac{1}{14+v}\right)\mathrm{d}v &= \frac{t}{20}+c \\
\frac{1}{28}\left(\ln(14+v)-\ln(14-v)\right) &= \frac{t}{20}+c
\end{align*}

Using \(v=7\) when \(t=0\):
\[
c=\frac{1}{28}\ln\left(\frac{21}{7}\right)=\frac{\ln 3}{28}
\]

Hence
\begin{align*}
\ln\left(\frac{14+v}{14-v}\right) &= \frac{7t}{5}+\ln 3 \\
\frac{14+v}{14-v} &= 3e^{7t/5} \\
14+v &= 42e^{7t/5}-3ve^{7t/5} \\
v\left(1+3e^{7t/5}\right) &= 14\left(3e^{7t/5}-1\right) \\
v &= \frac{14\left(3e^{7t/5}-1\right)}{3e^{7t/5}+1}
\end{align*}

\item Using \(v=\frac{\mathrm{d}x}{\mathrm{d}t}\), rewrite the result from part (a), whose denominator \(3e^{7t/5}+1\) is positive, to recognise a logarithmic integral:
\begin{align*}
\frac{\mathrm{d}x}{\mathrm{d}t}
&= -14+\frac{84e^{7t/5}}{3e^{7t/5}+1} \\
&= -14+20\frac{\frac{21}{5}e^{7t/5}}{3e^{7t/5}+1} \\
x &= -14t+20\ln\left(3e^{7t/5}+1\right)+d
\end{align*}

Since \(x=0\) when \(t=0\), \(d=-20\ln 4\). Therefore
\[
x=20\ln\left(\frac{3e^{7t/5}+1}{4}\right)-14t
\]

\item Using Newton's second law, with acceleration \(a=v\frac{\mathrm{d}v}{\mathrm{d}x}\):
\[
v\frac{\mathrm{d}v}{\mathrm{d}x}=\frac{196-v^2}{20}
\]

Separating the variables and integrating, with \(196-v^2>0\):
\begin{align*}
\int 1\,\mathrm{d}x &= \int\frac{20v}{196-v^2}\,\mathrm{d}v \\
x &= -10\ln(196-v^2)+h
\end{align*}

At entry, \(x=0\) and \(v=7\), so \(h=10\ln147\). Thus, when \(v=\frac{21}{2}\), the distance descended is
\[
x=10\ln\left(\frac{147}{196-\left(\frac{21}{2}\right)^2}\right)=10\ln\left(\frac{12}{7}\right)\,\mathrm{m}
\]
\end{enumerate}
{\raggedright\par\medskip\noindent\textbf{Final answers with \(g = 10\):} (a) \(v=10\sqrt{2}\,\frac{(10\sqrt{2}+7)e^{\sqrt{2}t}-(10\sqrt{2}-7)}{(10\sqrt{2}+7)e^{\sqrt{2}t}+(10\sqrt{2}-7)}\);\quad (b)~\(x=20\ln\left(\frac{(10\sqrt{2}+7)e^{\sqrt{2}t}+10\sqrt{2}-7}{20\sqrt{2}}\right)-10\sqrt{2}\,t\);\quad (c) \(10\ln\left(\frac{604}{359}\right)\,\mathrm{m}\)\par}
\end{djsSolution}

\newpage

% ---- Question 12 ----
\questionitem
% [Source: _QBT___Solns__Motion_Under_a_Variable_Force, Q12]

A package of mass $m\,\mathrm{kg}$ is released from rest and falls vertically. When its downward speed is $v\,\mathrm{m\,s^{-1}}$, it experiences a resistive force of magnitude $mkv^2\,\mathrm{N}$, where $k$ is a positive constant. Take the downward direction as positive. \\

\begin{enumerate}[label=\textbf{(\alph*)}]
  \item Show that, at time $t\,\mathrm{s}$ after release,
  \[
  v=\sqrt{\frac{g}{k}}\left(\frac{e^{2\sqrt{gk}\,t}-1}{e^{2\sqrt{gk}\,t}+1}\right)
  \]
  \marks[-40pt]{5}
\end{enumerate}
It is given that $k=0.0125\,\mathrm{m^{-1}}$. \\
\begin{enumerate}[label=\textbf{(\alph*)}, resume]
  \item Find the distance fallen by the package during the first $2$ seconds after its release. \marks{4} \\
\end{enumerate}

\vspace*{10pt}

\begin{djsSolution}
\begin{enumerate}[label=\textbf{(\alph*)}, start=1]
\item Using Newton's second law, with acceleration \(a=\frac{\mathrm{d}v}{\mathrm{d}t}\):
\begin{align*}
m\frac{\mathrm{d}v}{\mathrm{d}t}&=mg-mkv^2\\
\frac{\mathrm{d}v}{\mathrm{d}t}&=g-kv^2
\end{align*}

Let \(V=\sqrt{g/k}\) be the terminal speed. The package accelerates from rest towards this speed, so \(0\leq v<V\), and \(V+v\) and \(V-v\) are both positive.

Separating the variables and using partial fractions:
\begin{align*}
\int\frac{1}{V^2-v^2}\,\mathrm{d}v&=\int k\,\mathrm{d}t\\
\frac{1}{2V}\int\left(\frac{1}{V+v}+\frac{1}{V-v}\right)\,\mathrm{d}v&=\int k\,\mathrm{d}t\\
\frac{1}{2V}\ln\left(\frac{V+v}{V-v}\right)&=kt+c
\end{align*}

At release, \(t=0\) and \(v=0\), giving \(c=0\). Hence
\begin{align*}
\ln\left(\frac{V+v}{V-v}\right)&=2kVt\\
V+v&=e^{2kVt}(V-v)\\
v\left(1+e^{2kVt}\right)&=V\left(e^{2kVt}-1\right)\\
v&=\sqrt{\frac{g}{k}}\left(\frac{e^{2\sqrt{gk}\,t}-1}{e^{2\sqrt{gk}\,t}+1}\right)
\end{align*}

\item Let \(s\) be the distance fallen after \(t\) seconds. Using \(g=9.8\) and \(k=0.0125\) in part (a):
\begin{align*}
\frac{\mathrm{d}s}{\mathrm{d}t}=v
&=28\left(\frac{e^{0.7t}-1}{e^{0.7t}+1}\right)\\
&=\frac{56e^{0.7t}}{e^{0.7t}+1}-28
\end{align*}

The denominator \(e^{0.7t}+1\) is positive. Integrating by recognising the derivative of the denominator:
\begin{align*}
s&=\int\left(\frac{56e^{0.7t}}{e^{0.7t}+1}-28\right)\,\mathrm{d}t\\
&=80\ln\left(e^{0.7t}+1\right)-28t+d
\end{align*}

Using \(s=0\) when \(t=0\) gives \(d=-80\ln 2\), so
\[
s=80\ln\left(\frac{e^{0.7t}+1}{2}\right)-28t
\]

The distance fallen in the first \(2\) seconds is therefore
\[
s(2)=\left(80\ln\left(\frac{e^{1.4}+1}{2}\right)-56\right)\,\mathrm{m}
\approx18.2\,\mathrm{m}
\]
\end{enumerate}
\par\medskip\noindent\textbf{Final answers with \(g = 10\):} (b) \(\left(80\ln\left(\frac{e^{\sqrt{2}}+1}{2}\right)-40\sqrt{2}\right)\,\mathrm{m}\approx18.5\,\mathrm{m}\)
\end{djsSolution}

\newpage

% ---- Question 13 ----
\questionitem
% [Source: _QBT___Solns__Motion_Under_a_Variable_Force, Q13]

A particle $P$ of mass $16\,\mathrm{kg}$ moves in a straight line on a smooth horizontal plane. \\ 

At time $t=0$, the particle passes through a fixed point $O$ with velocity $\frac{5\sqrt{3}}{4}\,\mathrm{m\,s^{-1}}$. \\

From this instant, a horizontal force opposing the motion acts on $P$. The magnitude of the force is proportional to $\sqrt{25-4v^2}$, where $v\,\mathrm{m\,s^{-1}}$ is the velocity of $P$.  \\

When $v=2$, the magnitude of the force is $6\,\mathrm{N}$. The force ceases to act when the velocity has decreased to $\frac{5}{4}\,\mathrm{m\,s^{-1}}$. \\

\begin{enumerate}[label=\textbf{(\alph*)}]
  \item Show that, while the force is acting, the motion of $P$ is modelled by the differential equation
  \[
  \frac{1}{\sqrt{25-4v^2}}\frac{\mathrm{d}v}{\mathrm{d}t}=-\frac{1}{8}
  \]
  \marks[-30pt]{3}
  \item Hence find $v$ in terms of $t$, showing that your expression is valid for
  \[
  0\leq t\leq \frac{2\pi}{3}
  \]
  \marks[-20pt]{5}
  \item Determine the total distance travelled by $P$ during the interval $0\leq t\leq 8$. Give your answer in an exact form. \marks{6} \\
\end{enumerate}

\vspace*{10pt}

\begin{djsSolution}
\begin{enumerate}[label=\textbf{(\alph*)}, start=1]
  \item Take the direction of motion as positive. Let the magnitude of the opposing force be \(F=k\sqrt{25-4v^2}\), where \(k>0\).

  Using the given force when \(v=2\):
  \begin{align*}
  6&=k\sqrt{25-4(2)^2}\\
  6&=3k\\
  k&=2
  \end{align*}

  Using Newton's second law, with acceleration \(a=\frac{\mathrm{d}v}{\mathrm{d}t}\):
  \begin{align*}
  16\frac{\mathrm{d}v}{\mathrm{d}t}&=-2\sqrt{25-4v^2}\\
  \frac{1}{\sqrt{25-4v^2}}\frac{\mathrm{d}v}{\mathrm{d}t}&=-\frac18
  \end{align*}

  \item The standard result is
  \[
  \int\frac{1}{\sqrt{b^2-v^2}}\,\mathrm{d}v=\arcsin\left(\frac{v}{b}\right)+c\qquad(|v|<b)
  \]

  Separating the variables, factoring out the numerical scale and using it with \(b=\frac52\):
  \begin{align*}
  \frac{\mathrm{d}v}{\sqrt{25-4v^2}}&=-\frac18\,\mathrm{d}t\\
  \frac12\int\frac{1}{\sqrt{\left(\frac52\right)^2-v^2}}\,\mathrm{d}v
  &=-\frac18\int 1\,\mathrm{d}t\\
  \frac12\arcsin\left(\frac{2v}{5}\right)&=-\frac{t}{8}+c
  \end{align*}

  Substituting \(t=0\) and \(v=\frac{5\sqrt3}{4}\):
  \[
  c=\frac12\arcsin\left(\frac{\sqrt3}{2}\right)=\frac{\pi}{6}
  \]

  Hence
  \[
  \arcsin\left(\frac{2v}{5}\right)=\frac{\pi}{3}-\frac{t}{4}
  \]

  The force ceases when \(v=\frac54\), so
  \begin{align*}
  \arcsin\left(\frac12\right)&=\frac{\pi}{3}-\frac{t}{4}\\
  \frac{\pi}{6}&=\frac{\pi}{3}-\frac{t}{4}\\
  t&=\frac{2\pi}{3}
  \end{align*}

  Taking sine gives the required velocity while the force acts:
  \[
  v=\frac52\sin\left(\frac{\pi}{3}-\frac{t}{4}\right),
  \qquad 0\leq t\leq\frac{2\pi}{3}
  \]

  \item Let \(x\) be the displacement from \(O\) in the positive direction. Using \(v=\frac{\mathrm{d}x}{\mathrm{d}t}\) and integrating while the force acts:
  \begin{align*}
  \frac{\mathrm{d}x}{\mathrm{d}t}&=\frac52\sin\left(\frac{\pi}{3}-\frac{t}{4}\right)\\
  x&=10\cos\left(\frac{\pi}{3}-\frac{t}{4}\right)+d
  \end{align*}

  Substituting \(x=0\) at \(t=0\):
  \[
  0=10\cos\frac{\pi}{3}+d
  \quad\Longrightarrow\quad d=-5
  \]

  The velocity remains positive, so the distance travelled while the force acts is
  \[
  x\left(\frac{2\pi}{3}\right)=10\cos\frac{\pi}{6}-5=5\sqrt3-5
  \]

  After the force ceases, there is no horizontal force, so the particle continues at the constant velocity \(\frac54\,\mathrm{m\,s^{-1}}\). Therefore
  \begin{align*}
  \text{Total distance}&=(5\sqrt3-5)+\frac54\left(8-\frac{2\pi}{3}\right)\\
  &=\left(5\sqrt3+5-\frac{5\pi}{6}\right)\,\mathrm{m}
  \end{align*}
\end{enumerate}
\end{djsSolution}

\newpage

% ---- Question 14 ----
\questionitem
% [Source: _QBT___Solns__Motion_Under_a_Variable_Force, Q14]

A vehicle of mass $m\,\mathrm{kg}$ moves along a horizontal straight road. Its engine works at a constant rate of $2kmg\,\mathrm{W}$, where $k$ is a positive constant with units of speed. When the vehicle has speed $v\,\mathrm{m\,s^{-1}}$, it experiences a resistance of magnitude $\frac{mgv}{k}\,\mathrm{N}$. \\

Initially, the vehicle is at rest at $x=0$, where $x\,\mathrm{m}$ is the distance travelled. \\

\begin{enumerate}[label=\textbf{(\alph*)}]
  \item Show that
  \[
  v^2\dfrac{\mathrm{d}v}{\mathrm{d}x}=\dfrac{g}{k}\left(2k^2-v^2\right)
  \]
  \marks[-30pt]{3}
  \item Verify that
  \[
  gx=-kv+\dfrac{k^2}{\sqrt{2}}\ln\left(\dfrac{\sqrt{2}k+v}{\sqrt{2}k-v}\right)
  \]
  \marks[-40pt]{5}
  \item By using the work--energy principle, show that the time taken for the vehicle to accelerate from rest to a speed $V$, where $0<V<\sqrt{2}k$, is
  \[
  \dfrac{k}{2g}\ln\left(\dfrac{2k^2}{2k^2-V^2}\right)
  \]
  \marks[-40pt]{8}
\end{enumerate}

\vspace*{10pt}

\begin{djsSolution}
\begin{enumerate}[label=\textbf{(\alph*)}, start=1]
\item Take the direction of motion as positive. Using \(P=Fv\), the driving force is
\[
F=\frac{2kmg}{v}
\]

Using Newton's second law, with acceleration \(a=v\frac{\mathrm{d}v}{\mathrm{d}x}\):
\begin{align*}
mv\frac{\mathrm{d}v}{\mathrm{d}x}&=\frac{2kmg}{v}-\frac{mgv}{k}\\
v\frac{\mathrm{d}v}{\mathrm{d}x}&=\frac{g}{kv}\left(2k^2-v^2\right)\\
v^2\frac{\mathrm{d}v}{\mathrm{d}x}&=\frac{g}{k}\left(2k^2-v^2\right)
\end{align*}

\item Separating the variables in part (a):
\[
g\,\mathrm{d}x=k\frac{v^2}{2k^2-v^2}\,\mathrm{d}v
\]

Use algebraic division followed by partial fractions:
\begin{align*}
\frac{v^2}{2k^2-v^2}
&=-1+\frac{2k^2}{2k^2-v^2}\\
&=-1+\frac{k}{\sqrt{2}}\left(\frac{1}{\sqrt{2}k-v}+\frac{1}{\sqrt{2}k+v}\right)
\end{align*}

The resultant force is zero at the terminal speed \(\sqrt{2}k\). During acceleration from rest, \(0\leq v<\sqrt{2}k\), so \(\sqrt{2}k+v\) and \(\sqrt{2}k-v\) are both positive. Integrating gives
\begin{align*}
gx
&=k\int\left[-1+\frac{k}{\sqrt{2}}\left(\frac{1}{\sqrt{2}k-v}+\frac{1}{\sqrt{2}k+v}\right)\right]\,\mathrm{d}v\\
&=-kv+\frac{k^2}{\sqrt{2}}\left(\ln(\sqrt{2}k+v)-\ln(\sqrt{2}k-v)\right)+c
\end{align*}

When \(x=0\), \(v=0\), giving
\[
0=\frac{k^2}{\sqrt{2}}\left(\ln(\sqrt{2}k)-\ln(\sqrt{2}k)\right)+c
\quad\Longrightarrow\quad c=0
\]

Combining the logarithms:
\[
gx=-kv+\frac{k^2}{\sqrt{2}}\ln\left(\frac{\sqrt{2}k+v}{\sqrt{2}k-v}\right)
\]

\item Let \(W(v)\) be the work done by resistance from rest until the speed is \(v\). This work is negative because resistance opposes the motion. From part (a):
\[
\mathrm{d}x=\frac{kv^2}{g(2k^2-v^2)}\,\mathrm{d}v
\]

Hence
\begin{align*}
\mathrm{d}W&=-\frac{mgv}{k}\,\mathrm{d}x\\
\frac{\mathrm{d}W}{\mathrm{d}v}
&=-\frac{mv^3}{2k^2-v^2}
=mv-\frac{2mk^2v}{2k^2-v^2}
\end{align*}

Since \(0\leq v\leq V<\sqrt{2}k\), \(2k^2-v^2>0\). Integrating by recognising the derivative of \(2k^2-v^2\):
\[
W(v)=\frac12mv^2+mk^2\ln(2k^2-v^2)+d
\]

Using \(W(0)=0\) gives \(d=-mk^2\ln(2k^2)\), so
\[
W(V)=\frac12mV^2-mk^2\ln\left(\frac{2k^2}{2k^2-V^2}\right)
\]

The work--energy principle, with work positive when a force acts in the direction of motion, is
\[
KE_0+GPE_0+\sum W_{\text{other}}=KE_1+GPE_1
\]

Let the required time be \(t\). The engine does work \(Pt=2kmg\,t\). There is no change in gravitational potential energy. Substituting:
\begin{align*}
0+2kmg\,t+W(V)&=\frac12mV^2\\
2kmg\,t+\frac12mV^2-mk^2\ln\left(\frac{2k^2}{2k^2-V^2}\right)&=\frac12mV^2\\
2kmg\,t&=mk^2\ln\left(\frac{2k^2}{2k^2-V^2}\right)\\
t&=\frac{k}{2g}\ln\left(\frac{2k^2}{2k^2-V^2}\right)
\end{align*}
\end{enumerate}
\end{djsSolution}

\end{enumerate}
\end{document}
